%! Polynomial Functions and Complex Zeros — Unit 2, Lesson 2.4 — Guided Notes (Answer Key)
\documentclass[10pt]{article}
\usepackage{precalculus-article}
\usepackage{precalculus-key}   % pulls in -boxes; do NOT also load -boxes
\newcommand{\TallMath}[1]{$\displaystyle #1\rule[-1.4em]{0pt}{3.2em}$}

\begin{document}

\pageheader{Unit 2, Lesson 2.4}{Guided Notes --- Answer Key}

\vspace{0.12in}

\begin{objectivebox}
By the end of this lesson, I will be able to\ldots
\begin{enumerate}[itemsep=2pt]
  \item Explain why a polynomial can have zeros that never show up as
        \ans{$x$-intercepts} on its graph.
  \item Use the Fundamental Theorem of Algebra to count the \ans{$n$}
        complex zeros of a polynomial of degree $n$.
  \item Use \ans{conjugate} pairs to name a missing zero and to build a
        polynomial from a list of zeros.
  \item Decide whether a polynomial function is \ans{even},
        \ans{odd}, or neither.
\end{enumerate}
\end{objectivebox}

\vspace{0.1in}

\boxguard
\begin{vocabbox}
Fill in each term as we define it together.
\par\vspace{2pt}
\noindent\textbf{\textcolor{plum}{Imaginary unit:}}\\[1pt]
\ansline{The number $i$, defined by $i^2 = -1$.}\\[3pt]
\noindent\textbf{\textcolor{plum}{Complex number:}}\\[1pt]
\ansline{Any number $a + bi$ with $a$, $b$ real; $a$ is the real part, $b$ the imaginary part.}\\[3pt]
\noindent\textbf{\textcolor{plum}{Non-real zero:}}\\[1pt]
\ansline{A zero $a + bi$ with $b \ne 0$ --- not an $x$-intercept, and not on the graph at all.}\\[3pt]
\noindent\textbf{\textcolor{plum}{Complex conjugate:}}\\[1pt]
\ansline{The partner $a - bi$ of $a + bi$; only the sign in front of the $i$ changes.}\\[3pt]
\noindent\textbf{\textcolor{plum}{Fundamental Theorem of Algebra:}}\\[1pt]
\ansline{A polynomial of degree $n$ has exactly $n$ complex zeros, counting multiplicities.}\\[3pt]
\noindent\textbf{\textcolor{plum}{Even function:}}\\[1pt]
\ansline{$f(-x) = f(x)$ for every $x$; the graph is symmetric over the line $x = 0$.}\\[3pt]
\noindent\textbf{\textcolor{plum}{Odd function:}}\\[1pt]
\ansline{$f(-x) = -f(x)$ for every $x$; the graph is symmetric about the point $(0,0)$.}
\end{vocabbox}

\vspace{0.1in}

\boxguard[34]
\begin{hookbox}
Below is the graph of $p(x) = x^4 - 3x^2 - 4$, a polynomial of degree 4. It
crosses the $x$-axis exactly twice.

\begin{center}
\begin{tikzpicture}[x=1.0cm, y=0.28cm]
  \draw[linegray, very thin, xstep=1, ystep=2] (-2.4,-7) grid (2.4,5);
  \draw[->, thick] (-2.7,0) -- (2.8,0) node[below right] {\small $x$};
  \draw[->, thick] (0,-7.6) -- (0,5.8) node[above] {\small $p(x)$};
  \foreach \x in {-2,-1,1,2} \node[below] at (\x,-0.3) {\small \x};
  \foreach \y in {-6,-4,-2,2,4} \node[left] at (0,\y) {\small \y};
  \draw[very thick, plum] plot [smooth] coordinates
    {(-2.2,4.9) (-2,0) (-1.6,-5.13) (-1.2,-6.25) (-0.8,-5.51)
     (-0.4,-4.45) (0,-4) (0.4,-4.45) (0.8,-5.51) (1.2,-6.25)
     (1.6,-5.13) (2,0) (2.2,4.9)};
  \fill[plum] (-2,0) circle (2.5pt);
  \fill[plum] (2,0) circle (2.5pt);
\end{tikzpicture}
\end{center}

Lesson 2.3 tied every zero to a factor. A degree-4 polynomial has four factors'
worth of zeros --- but the graph shows only two. Where are the other two?

\ansline{They are there, but they are not real numbers: $p(x) = (x-2)(x+2)(x^2+1)$, and}
\ansline{$x^2 + 1 = 0$ gives $x = i$ and $x = -i$. A graph can only show real zeros.}
\end{hookbox}

\vspace{0.1in}

\boxguard[30]
\begin{notesbox}{1. A Complex-Number Refresher}
From Algebra 2: the \textbf{imaginary unit} $i$ is the number with
$i^2 = $ \ans{$-1$}. A \textbf{complex number} has the form $a + bi$, where
$a$ and $b$ are real numbers.

\begin{itemize}[itemsep=4pt]
  \item $a$ is the \ans{real} part and $b$ is the \ans{imaginary} part.
  \item Every real number is already complex: $7 = 7 + 0i$. A number is
        \textbf{non-real} exactly when $b$ \ans{$\ne$} $0$.
  \item The \textbf{complex conjugate} of $a + bi$ is \ans{$a - bi$}. Only the
        sign in front of the $i$ changes.
\end{itemize}

\textbf{Example 1.} Show that $i$ is a zero of $p(x) = x^2 + 1$. Nothing new
here --- a zero is still an input that makes the output $0$.

\begin{work}
  p(i) &= (i)^2 + 1 \\
       &= -1 + 1 \\
       &= 0
\end{work}

\textbf{Example 2.} Multiply the conjugate pair $3 + 2i$ and $3 - 2i$.

\begin{work}
  (3+2i)(3-2i) &= 9 - 6i + 6i - 4i^2 \\
               &= 9 - 4(-1) \\
               &= 13
\end{work}

The $i$ terms cancelled and the answer came out \ans{real}. Remember this
--- it is the reason for Section 3.
\end{notesbox}

\vspace{0.1in}

\boxguard[30]
\begin{notesbox}{2. The Fundamental Theorem of Algebra --- Exactly $n$ Zeros}
\textbf{Fundamental Theorem of Algebra.} A polynomial function of degree $n$
has \textbf{exactly} \ans{$n$} complex zeros, counting multiplicities.

\begin{itemize}[itemsep=4pt]
  \item ``Exactly,'' not ``at most.'' The count is guaranteed once we allow
        non-real numbers as inputs.
  \item \textbf{Real} zeros are the ones you can see: they are the
        \ans{$x$-intercepts} of the graph.
  \item \textbf{Non-real} zeros are not on the graph \textit{at all}. A graph
        can never show them, so they must be found from the factors.
\end{itemize}

\textbf{The hook, resolved.} $p(x) = x^4 - 3x^2 - 4 = (x-2)(x+2)(x^2+1)$.
Setting $x^2 + 1 = 0$ gives $x^2 = -1$, so $x = $ \ans{$\pm i$}. All four
zeros: $2$, $-2$, and the pair \ans{$i$, $-i$}.

\medskip
\textbf{Zero census.} Complete the table. Count real zeros \textit{with
multiplicity}, then subtract from the degree.

\smallskip
\renewcommand{\arraystretch}{1.5}
\begin{center}
\begin{tabular}{|l|c|c|c|c|}
\hline
\textbf{Factored form} & \textbf{Degree} & \textbf{\# real} & \textbf{\# non-real} & \textbf{Total} \\
\hline
$(x-2)(x+2)(x^2+1)$ & \ans{4} & \ans{2} & \ans{2} & \ans{4} \\
\hline
$(x+1)^3(x^2+4)$ & \ans{5} & \ans{3} & \ans{2} & \ans{5} \\
\hline
$x(x^2+9)$ & \ans{3} & \ans{1} & \ans{2} & \ans{3} \\
\hline
$(x-1)^2(x^2+1)^2$ & \ans{6} & \ans{2} & \ans{4} & \ans{6} \\
\hline
\end{tabular}
\end{center}
\smallskip

\textbf{The move to memorize:} degree $-$ (real zeros with multiplicity) $=$
number of \ans{non-real} zeros.
\end{notesbox}

\vspace{0.1in}

\boxguard[30]
\begin{notesbox}{3. Non-Real Zeros Come in Pairs}
\textbf{The conjugate pair rule.} If $a + bi$ is a zero of a polynomial with
real coefficients, then \ans{$a - bi$} is also a zero.

\begin{itemize}[itemsep=4pt]
  \item Consequence 1: the number of non-real zeros of such a polynomial is
        always \ans{even} --- they cannot arrive one at a time.
  \item Consequence 2: a polynomial of \textbf{odd} degree has an odd total
        number of zeros, so they cannot all be non-real. It must have at least
        one \ans{real} zero --- its graph \textit{has} to cross the
        $x$-axis.
  \item Look back at Section 1: a conjugate pair multiplies out to a factor
        with real coefficients, which is why the pair is forced.
\end{itemize}

\textbf{Example 3.} A polynomial $p$ has real coefficients, degree 3, leading
coefficient 1, and zeros $x = 2$ and $x = 1 + 3i$. Name the third zero, then
write $p(x)$.

Third zero: \ans{$1 - 3i$}

\begin{work}
  p(x) &= (x-2)\big(x - (1+3i)\big)\big(x - (1-3i)\big) \\
       &= (x-2)\big((x-1) - 3i\big)\big((x-1) + 3i\big) \\
       &= (x-2)\big((x-1)^2 - 9i^2\big) \\
       &= (x-2)\big(x^2 - 2x + 1 + 9\big) \\
       &= (x-2)(x^2 - 2x + 10) \\
       &= x^3 - 4x^2 + 14x - 20
\end{work}

Every coefficient came out \ans{real} --- exactly what the conjugate pair
guarantees.
\end{notesbox}

\vspace{0.1in}

\boxguard[34]
\begin{notesbox}{4. Even and Odd Polynomial Functions}
\begin{itemize}[itemsep=4pt]
  \item $f$ is \textbf{even} when $f(-x) = $ \ans{$f(x)$} for every $x$. Its
        graph is symmetric over the line \ans{$x = 0$}.
  \item $f$ is \textbf{odd} when $f(-x) = $ \ans{$-f(x)$} for every $x$. Its
        graph is symmetric about the point \ans{$(0,0)$}.
  \item Most functions are \textbf{neither}. ``Odd'' is not the same as ``not
        even.''
\end{itemize}

\textbf{Example 4.} Test $g(x) = x^3 - 4x$ by substituting $-x$.

\begin{work}
  g(-x) &= (-x)^3 - 4(-x) \\
        &= -x^3 + 4x \\
        &= -(x^3 - 4x) \\
        &= -g(x)
\end{work}

So $g$ is \ans{odd}.

\medskip
\textbf{The polynomial shortcut.} Look only at the exponents:

\begin{itemize}[itemsep=4pt]
  \item Every exponent \textbf{even} $\Rightarrow$ the function is
        \ans{even}. \quad (Careful: a constant term is $x^0$, an
        \textit{even} power.)
  \item Every exponent \textbf{odd} $\Rightarrow$ the function is
        \ans{odd}.
  \item A \textbf{mix} of even and odd exponents $\Rightarrow$ \ans{neither}.
\end{itemize}

\textbf{Reading the symmetry.} Graph A is $y = x^4 - 4x^2$; graph B is
$y = x^3 - 3x$.

\begin{center}
\begin{tikzpicture}[x=0.55cm, y=0.25cm]
  \draw[->, thick] (-2.6,0) -- (2.7,0) node[below right] {\small $x$};
  \draw[->, thick] (0,-4.8) -- (0,5.0) node[above] {\small $y$};
  \draw[very thick, plum] plot [smooth] coordinates
    {(-2.2,4.07) (-2,0) (-1.7,-3.21) (-1.4,-4.0) (-1,-3) (-0.5,-0.94)
     (0,0) (0.5,-0.94) (1,-3) (1.4,-4.0) (1.7,-3.21) (2,0) (2.2,4.07)};
  \node[below] at (0,-5.4) {\small \textbf{Graph A}};
\end{tikzpicture}
\hspace{1.1cm}
\begin{tikzpicture}[x=0.55cm, y=0.25cm]
  \draw[->, thick] (-2.6,0) -- (2.7,0) node[below right] {\small $x$};
  \draw[->, thick] (0,-4.8) -- (0,5.0) node[above] {\small $y$};
  \draw[very thick, plum] plot [smooth] coordinates
    {(-2,-2) (-1.8,-0.43) (-1.5,1.13) (-1,2) (-0.5,1.38) (0,0)
     (0.5,-1.38) (1,-2) (1.5,-1.13) (1.8,0.43) (2,2)};
  \node[below] at (0,-5.4) {\small \textbf{Graph B}};
\end{tikzpicture}
\end{center}

Graph A is \ans{even} (mirror over the $y$-axis); graph B is
\ans{odd} (half-turn about the origin).
\end{notesbox}

\vspace{0.1in}

\boxguard
\begin{practicebox}
\begin{enumerate}[itemsep=8pt]
  \item Let $p(x) = (x-6)(x^2 + 2x + 5)$. Give the degree, then list all
        complex zeros.

        \begin{work}
          x^2 + 2x + 5 &= 0 \\
                     x &= \frac{-2 \pm \sqrt{4 - 20}}{2} \\
                       &= \frac{-2 \pm 4i}{2} \\
                       &= -1 \pm 2i
        \end{work}

        Degree $= $ \ans{3} \qquad Zeros: \ans{$6$, $-1 + 2i$, $-1 - 2i$}

  \item Classify each polynomial as \textbf{even}, \textbf{odd}, or
        \textbf{neither}.

        \begin{enumerate}[label=(\alph*), itemsep=4pt]
          \item $f(x) = x^4 - 7x^2 + 1$ \quad \ans{even}
          \item $g(x) = x^5 - x$ \quad \ans{odd}
          \item $h(x) = x^4 + x$ \quad \ans{neither}
        \end{enumerate}
\end{enumerate}
\end{practicebox}

\end{document}
